\section{Sequences with repetition}

A sequence with repetition is an ordered list in which each position is filled separately from the same set of symbols.

\begin{definitionbox}
Let \(S\) be a set with \(n\) elements. A \textbf{sequence of length \(k\) with repetition} over \(S\) is an ordered \(k\)-tuple
\[
(a_1,a_2,\ldots,a_k),
\]
where each \(a_i\in S\). The same element may occur many times.
\end{definitionbox}

The word ``ordered'' means that positions matter. The first position, the second position, and the third position are different places in the description. The phrase ``with repetition'' means that choosing one symbol at one position does not remove it from later positions.

Now we derive the count. To construct one sequence,
\begin{itemize}
    \item choose \(a_1\): there are \(n\) choices;
    \item choose \(a_2\): there are again \(n\) choices, because repetition is allowed;
    \item continue in the same way up to \(a_k\): there are \(n\) choices at each position.
\end{itemize}
By the multiplication principle, the number of complete sequences is
\[
\underbrace{n\cdot n\cdots n}_{k\text{ factors}}=n^k.
\]
Thus the formula \(n^k\) is not a separate rule to memorize. It is the multiplication principle applied to \(k\) ordered positions, each with the same \(n\) choices.

This model is used when:
\begin{itemize}
    \item order is recorded,
    \item each position is chosen separately,
    \item repetition is allowed,
    \item the outcome is a full sequence.
\end{itemize}

Typical examples include PIN codes, repeated coin tosses, repeated die rolls, strings over an alphabet, and repeated selections with replacement.

\begin{examplebox}
The space of possibilities for three tosses of a coin is
\[
\{H,T\}^3.
\]
It has
\[
2^3=8
\]
elementary outcomes. Each outcome is a sequence, for example
\[
(H,T,H).
\]
The outcomes \((H,T,H)\) and \((T,H,H)\) are different because the order of positions is recorded.

A tree representation shows the same count visually. At the first toss there are \(2\) branches, from each of these there are \(2\) branches for the second toss, and from each of those there are \(2\) branches for the third toss. The tree has
\[
2\cdot 2\cdot 2=8
\]
terminal branches, one for each complete sequence.
\end{examplebox}

\section{Ordered selections without repetition}

Sometimes only some objects are chosen, but the order or position of the chosen objects matters.

\begin{definitionbox}
An \textbf{ordered selection without repetition} of length \(k\) from \(n\) distinct objects is an ordered list of \(k\) different objects chosen from the \(n\) available objects.
\end{definitionbox}

To derive the count, construct the ordered list position by position. There are \(n\) choices for the first position. Once the first object has been used, it cannot be used again, so there are \(n-1\) choices for the second position. After two objects have been used, there are \(n-2\) choices for the third position. Continuing in this way, the \(k\)-th position has
\[
n-(k-1)=n-k+1
\]
choices.
Therefore the number of ordered selections without repetition is
\[
P(n,k)=n(n-1)(n-2)\cdots(n-k+1).
\]
This product can also be written using factorials. Since
\[
n!=n(n-1)\cdots(n-k+1)(n-k)(n-k-1)\cdots 1,
\]
dividing by \((n-k)!\) removes the unused tail:
\[
P(n,k)=\frac{n!}{(n-k)!}.
\]

This model is used when:
\begin{itemize}
    \item only some objects are selected,
    \item order or position matters,
    \item no object may be used twice.
\end{itemize}

\begin{examplebox}
Gold, silver, and bronze medals are assigned among \(12\) runners. An outcome is not a set of three runners. It is an ordered triple
\[
(\text{gold},\text{silver},\text{bronze}).
\]
There are \(12\) choices for gold, then \(11\) choices for silver, then \(10\) choices for bronze. Hence the number of possible assignments is
\[
12\cdot 11\cdot 10=P(12,3).
\]
\end{examplebox}

In some European terminology, ordered selections are called variations. In these notes we use the name ordered selections without repetition, or \(k\)-permutations, because it describes directly what the outcome is.

\section{Arrangements of all objects}

A permutation is an arrangement of all objects in a row.

\begin{definitionbox}
A \textbf{permutation} of \(n\) distinct objects is a linear arrangement of all \(n\) objects.
\end{definitionbox}

This is the special case of an ordered selection without repetition in which all objects are used. Thus \(k=n\), and
\[
P(n,n)=n(n-1)(n-2)\cdots 2\cdot 1=n!.
\]
So there are
\[
n!
\]
permutations of \(n\) distinct objects.

\begin{examplebox}
If \(7\) students are arranged in a line, then one outcome is a complete ordered list of all students. The number of arrangements is
\[
7!.
\]
This is not because \(7!\) is a name attached to the word ``arrangement''. It is because there are \(7\) choices for the first position, \(6\) for the second, \(5\) for the third, and so on until all positions are filled.
\end{examplebox}

Permutations are used when:
\begin{itemize}
    \item all objects are used,
    \item order matters,
    \item the objects are distinguishable,
    \item the arrangement is linear.
\end{itemize}

\section{Circular arrangements}

A circular arrangement is different from a linear arrangement because rotating the whole arrangement does not create a new relative seating order.

\begin{examplebox}
Suppose four people \(A,B,C,D\) sit around a round table. The linear lists
\[
(A,B,C,D),\qquad (B,C,D,A),\qquad (C,D,A,B),\qquad (D,A,B,C)
\]
represent the same circular arrangement if only relative positions matter.
\end{examplebox}

For \(n\) distinct objects placed around a circle, with rotations identified, the number of circular arrangements is
\[
(n-1)!.
\]
One way to see this is to fix one object as a reference point. This does not reduce any real choice; it only chooses one representative from each group of rotated descriptions. Once the reference object is fixed, the remaining \(n-1\) objects can be placed around it in a linear order:
\[
(n-1)(n-2)\cdots 2\cdot 1=(n-1)!.
\]

\begin{remarkbox}
A circular arrangement requires a modelling decision. If seats are numbered or if there is a distinguished position, then rotations may become different outcomes and the linear count may be appropriate.
\end{remarkbox}
